\section{Fight over Origins}

\textbf{Julius Evola} opens the chapter The Arctic Myth from Il mito del sangue by describing the state of prehistoric research:

\begin{quotex}
Oswald Menghin, rector of the University of Vienna, wrote these characteristic words, ``More than any other discipline, the science of prehistory is being brought back, and still more must be brought back, to the center of the spiritual battle of our time. I don't believe I am mistaken in affirming that general prehistory will be the science that will guide the next generations." In the most recent years, in many circles a significant drive to return to the origins has taken place. The origins, in this case, appear under a special, spiritual light. We return to the intuition that, in primordial times, meaning and symbols existed in a still pure state that was then lost, obfuscated, or altered. Prehistoric research, moved from a plane of disanimiated scientistic-archeological or anthropological positivism to a plane of spiritual synthesis, promised therefore to open up new horizons for the true history of civilization. 

\end{quotex}
In a sense this is still true, as the battle for control over the narrative of prehistory rages on. However, the prehistory preferred by Evola, based on the Arctic myth of \textbf{Fabre d'Olivet}, \textbf{Hermann Wirth}, and \textbf{Bal Tilak}, has been relegated to the world of cranks and conspiracy theorists. Unfortunately, the alternative myths of origin, that have adherents today, remain far from any ``plane of spiritual synthesis".

The preferred ``scientific" myth of educated (read indoctrinated) people today involves the neo-gnostic belief in a big bang that led to an undifferentiated humanity by the random walk theorized by \textbf{Charles Darwin}. In this story, man arises from a great, but incomprehensible, cosmic accident, thrown into a world he cannot understand, and placed in unjust situation, both individually and socially. Hence, the task of man is to develop a higher consciousness that sees beyond the creator of these injustices into a future utopia of liberty, equality, and brotherhood.

The competing myth in the West sees the world as created just, but injustice then entered the world through one man's sin. There are different degrees of literalness that goes with this story, from the adamant to the flexible who try to accommodate the scientistic view. However, they all share the same materialist, empirical presumptions, so no reconciliation or synthesis on that level is possible. Evola points to the alternative:

\begin{quotex}
In support of the Nordic idea a new myth was therefore necessary, a myth as audacious, as it was rich, complex and articulate, was now the material of the know-how to dominate and organize according to a unique explicative principle. Such a myth was forged by the Dutchman Hermann Wirth with the recovery of the arctic theory … with the reconstruction of the origin, history and civilization of the Nordic-Atlantic race. The theory of Wirth is to be considered as a bold surprise, whose profound directive impulse referred to extrascientific intuitions which then sought to be justified through a very laborious philological, anthropo-geological, mythological, and symbological apparatus. 

\end{quotex}
While at one time, it may have had great explanatory power, it has little serious traction today. Tilak's so-called ``Aryan Invasion Theory" is now rejected by academics, no one reads Fabre d'Olivet, and Wirth's book is effectively unavailable. It is pointless to simply repeat what has been said and it is ludicrous to associate this with stories of UFOs or underground worlds. To be taken seriously, its alleged explanatory powers must be used to explain things.

First of all, we must not look East, but rather to the North. There is a primordial monotheistic tradition that took various forms in the period of migrations, just as languages did, while retaining similar structures. We have brought to light aspects of those traditions, emphasizing their similarities and continuity.

There is an unfortunate tendency in most of those interested in the concepts of tradition, and that is over-intellectualization. But the realm of the spiritual is inaccessible to this type of analysis. What is necessary is the imaginary reconstruction of spiritual states through a type of meditation, based on myths, symbols, rites, legends, and so on. The men capable of doing that will be the creators of the new world, not those addicted to the Demon of Dialectics\footnote{\url{https://gornahoor.net/?p=3053}}.



\flrightit{Posted on 2011-09-24 by Cologero }

\begin{center}* * *\end{center}

\begin{footnotesize}\begin{sffamily}



\texttt{logres on 2011-09-24 at 19:25 said: }

It's interesting that Lewis and Tolkien were looking ``north" as well – they are practically all that Christendom has left these days. Isn't there a northern myth in the East, as well?


\hfill

\texttt{Zmaj on 2011-09-25 at 04:15 said: }

I don't see the way in which you oppose Guenon by looking to the North. If he emphasized the East, it was as a bastion of remaining traditional forms, especially those that reflect the purity of primordiality in their initiatory sophistication. He explained his view that the primary movement of the tradition throughout our cyclical development was from the Northern home, to the West, and finally Eastward as Western degeneracy accelerated. So what do you look to in the North? Might one argue that you must first look to the East (or the esoteric West) for a vehicle to guide you North?

Guenon related knowledge of Central and South American peoples that made reference to an Atlantic home. As for a northern myth in the East, if we allow that the northern myth corresponds to the polar myth (even without reference to names like Atlantis and Borea), we find the symbolism of the mountain throughout the east from Islam to India to Japan, with Polaris holding great significance in Taoism. These symbols represent a place of supernatural origin in these traditions, just as Borea represents the earthly origin of the derived supernatural traditional forms.

Concerning the Southern Negroid races, will someone comment on their place in relation to the Northern myth? I understand that they are the decadent legacy of a previous cycle of developed civilization; what does this say about their condition during/relations with Borea and Atlantis and the related cataclysms? Was their earlier civilization a counterpart to Borea or one of its earliest derivatives? Guenon treated the topic of savages somewhat enigmatically, probably with good reason.

As for today's common beliefs, what do you all make of the article here?

\url{http://www.abc.net.au/news/2011-09-23/aboriginal-dna-dates-australian-arrival/2913010}


\hfill

\texttt{Matt on 2011-09-26 at 17:37 said: }

Is it accurate to say that the Aryan Invasion/migration theory is now rejected by academics? From my understanding, the vast majority of linguists still hold to it, and great bulk, possibly the majority of archaeologists and anthropologists also hold to it.


\hfill

\texttt{Cologero on 2011-09-26 at 23:10 said: }

You tell me, Matt. Do you know of any specific academic who has defended that idea. We can start with Amazon, and we don't see many supporters on the first page. Perhaps an academic may support the idea of an invasion from Central Asia, but other than Tilak, who has supported the idea of Hyperborean origin? It won't fit with the ``out of Africa" theory, which probably most academics implicitly accept despite a complete lack of evidence.

But the question I asked is why does Evola refer to it as a myth rather than a theory? And he claims it is the result of ``extrascientific intuitions", which is only afterwards justified with a ``very laborious philological, anthropo-geological, mythological, and symbological apparatus". This would seem to rule out academics.


\hfill

\texttt{Matt on 2011-09-27 at 19:41 said: }

Oh I thought you were referring to the migrations of the Indo-Aryans into northern India. Sorry I'm just so used to hearing people talk about that when ``Aryan Invasions" comes up that I thought thats what you were referring to it.

But yes, when it comes Hyperborea and academics, your statement is quite accurate.

[I'm still curious about a recent academic book on the AIT, wherever they may have originated.]


\hfill

\texttt{Matt on 2011-09-27 at 19:48 said: }

And as for the first page of Amazon, its not surprising that the books listed that present whats called the ``Out of India" theory are all written by what appears to be Indian nationalists (its been long known that that theory has been tightly connected with the nationalist movement).


\hfill

\texttt{Cologero on 2011-09-27 at 20:02 said: }

How long has that been known? Wasn't Bal Tilak a leader in the India independence movement?


\hfill

\texttt{Matt on 2011-09-28 at 01:00 said: }

Yes, Tilak was apart of the independence movement. To be honest, I'm not sure how well his view of where the Vedas were composed and to where they're composers came from resonates within the nationalist movement (though I bet it doesn't resonate strongly).

But getting back to Out of India, that theory did seem to arise right at the time when the Indian nationalist movement was blossoming, and therefore its ties with the movement have always stuck to it. Whats interesting th0ough it that a recent genetic study showed that the genetic marker R1a (which is associated with the Indo-Aryans and Indo-Europeans as a whole), seems to have originated in northern India, which would give some weight to the out of india theory (though it still has a significant amount of linguistic, historical, archaeological, and genetic evidence going against it). Interesting nonetheless. But as you and Evola have pointed out, none of these competing theories really get at any idea of a spiritual synthesis (and they most likely never will).

What was the recent academic book on the AIT that you mentioned?


\hfill

\texttt{Perennial on 2011-09-29 at 02:52 said: }

Does anyone have any thoughts as to how Adam and Eve would fit into a Northern/Aryan origin? I do not recall either Guenon or Evola addressing this specific issue. I am curious as to how Nordic idea would fit into a primordial fall.


\hfill

\texttt{Boreas on 2011-09-29 at 05:16 said: }

``Does anyone have any thoughts as to how Adam and Eve would fit into a Northern/Aryan origin? I do not recall either Guenon or Evola addressing this specific issue. I am curious as to how Nordic idea would fit into a primordial fall."

Are you asking that is a similar myth found in the Nordic traditions? Are you aware of the story of Ask and Embla (Adam and Eve) to whom Hroptri-Odhínn (``the cosmic Christ" / mystical christos) gave life, light and intelligence?

I read about the myth of Heimdall the last weekend and the myth and Heimdall's character resembled in my opinion a lot the myth of the fall of Lucifer (without the `war in heaven' part) and it also remainded me about the theosophical teachings of the humanity of the moon-cycle. (Heimdall is said to be originally a moon-god).


\hfill

\texttt{Eagle on 2011-10-02 at 17:00 said: }

[Nordics, such as Snorri Sturluson, wrote down the Norse myths. Cologero]

``Are you asking that is a similar myth found in the Nordic traditions? Are you aware of the story of Ask and Embla (Adam and Eve) to whom Hroptri-Odhínn (``the cosmic Christ" / mystical christos) gave life, light and intelligence?"

We have no way of knowing how much Christian alteration we have of the Norse myths. Consider who copied them down and when.


\hfill

\texttt{Boreas on 2011-10-03 at 04:23 said: }

``We have no way of knowing how much Christian alteration we have of the Norse myths. Consider who copied them down and when."

That's true, but we must stick to what we have of them, make the best out of it, and then make our own further examinations and studies. I'm not saying that Odin or Balder = Christ or Heimdall = Lucifer (prior to fall, which I consider to be a cosmic and not a ``moral" incidence), but there are a lot of convergences and similarities in the myths and symbols.


\hfill

\texttt{Charlotte Cowell on 2011-10-04 at 10:58 said: }

The idea that `civilised' history can only be charted from around 3,000 BC is nonsensical sd clearly – obviously – there were more ancient peoples than this around way past this point. For Atlantis it's fairly simple, you're looking at a zone that stretches from the Caribbean coast to the Senegal region, possibly a series of large islands. The Bay of Bengal area is another massive `hot spot' for mega-ancient civilisations and in fact sunken cities are cropping up all over the place, not to mention pyramids. The largest pyramids in Europe – gigantic things far bigger and even more astronomically precise than the Great Pyramid – were recently unearthed in Bosnia, and China has kept its 150 plus pyramids well hidden from view for reasons best known to China. Last I saw they were planting trees on them to hide them even more. The Sphinx, as we know, is around 10,000 years old and there's no indisputable evidence that the Great Pyramids were built for Cheops. They could be waaaaaay older than that, but most people will suck up the info. fed to them from the Egyptian authorities that a few thousand slaves built the things in 20 years with sand shelves, logs, a few chisels and big round stones. Big round stones to create the most mathematically precise sculptures and buildings ever seen on Earth. By the way, did you know the great pyramid actually has 8 sides not 4? It doesn't become clear until the Autumn equinox (for a few seconds), when one half of one side is in shadow and the other in light. It is even possible that we live in a binary star system that the ancients knew about and that we are yet to discover….(possible, I said, not certain yet!). Perhaps Sirius is more than we previously imagined. Joscelyn Godwin argues that Hyperborea was in ancient Britain, which does not seem unreasonable given that we have a large solar (Apollonian) monument at Stonehenge. Also regarding ancient destroyed worlds, in a certain part of India, near to Bangladesh, I think, there is evidence of an immense nuclear cataclysm many thousands of years ago, based on analysis of the soil and levels of radiation. The main problem for human beings in sorting out the order of things is that (for the most part) our memories just ain't what they used to be. It is likely that there have been past disasters of such magnitude – and so terrifying – that there is a collective mental block about them ingrained into our DNA….Agree that The Fall occurs cosmically but also within each microcosm, and there have been many reverberations throughout known and unknown history as successive generations encounter the same stumbling blocks.


\hfill

\texttt{Charlotte Cowell on 2011-10-04 at 17:16 said: }

Just an addendum about the twin sun theory: For more info from ancient sources read ``The Holy Science" by Sri Yukteswar.

Yuga Theory

Sri Yukteswar's introduction to The Holy Science includes his explanation of the Yuga Cycle – revolutionary because of his premise that the earth is now in the age of Dwapara Yuga, not the Kali Yuga that most Indian pundits believe to be the current age.[4] His theory is based on the idea that the sun ``takes some star for its dual and revolves round it in about 24,000 years of our earth – a celestial phenomenon which causes the backward movement of the equinoctial points around the zodiac."[1] The common explanation for this celestial phenomenon is precession, the `wobbling' rotating movement of the earth axis. Research into Sri Yukteswar's explanation is being conducted by the Binary Research Institute.


\hfill

\texttt{Boreas on 2011-10-06 at 03:04 said: }

Charlotte: As I've been studying the theory of the four yugas, I've myself come to a somewhat certain conclusion – with no authority on the matter of course – that there are yugas within yugas within yugas within yugas – all the way to the most minimal part of time cycles. What I mean is that the four time cycles repeat themselves in every yuga: there is the golden age of the Kali Yuga, the silver age of the Kali Yuga, the bronze age of the Kali Yuga, and the iron age of the Kali Yuga; and all these four cycles and the whole of the particular Kali Yuga can be within a greater Dwapara Yuga, for example. (I would like to hear thoughts about the exactness of this thought pattern from someone who is knowledgeable in these matters.)


\hfill


\end{sffamily}\end{footnotesize}
